\subsection{Verificação técnica do assistente}
\label{subsec:verificacao-interface}

Em 12 de setembro de 2026, foi verificada a rota local \texttt{/language-creator}, sem autenticação, em Chromium 153.0.8010.12, usando Playwright para interação e axe-core 4.13.0 para verificações automatizadas. Foram combinadas duas larguras (1440 e 390 pixels CSS) com os temas claro e escuro, em quatro cenários. O portal de ferramentas de desenvolvimento do Next.js foi ocultado e excluído da análise. A verificação abrangeu a etapa inicial do assistente e a navegação para entrada/saída; não incluiu contas, turmas, submissões ou leitor de tela.

Nos quatro cenários, cinco acionamentos de Tab alcançaram o botão da primeira etapa; Tab seguido de Enter ativou a segunda etapa. A lista de etapas apresentou direção vertical em 1440 pixels e horizontal em 390 pixels. A largura de rolagem do documento coincidiu com a largura da janela nos dois tamanhos. Esses resultados caracterizam o estado testado, não todos os componentes e estados da plataforma. A Figura~\ref{fig:responsividade-wizard} apresenta as capturas correspondentes no tema escuro.

\begin{figure}[p]
\centering
\begin{subfigure}{\textwidth}
\centering
\includegraphics[width=\textwidth,height=0.40\textheight,keepaspectratio]{wizard-desktop-atual.png}
\caption{Janela de 1440 por 1000 pixels CSS: etapas na lateral e conteúdo em colunas.}
\end{subfigure}

\medskip
\begin{subfigure}{\textwidth}
\centering
\includegraphics[height=0.32\textheight,keepaspectratio]{wizard-mobile-atual.png}
\caption{Janela de 390 por 844 pixels CSS: etapas em linha e cartões empilhados.}
\end{subfigure}
\caption{Adaptação do assistente na versão verificada em setembro de 2026.}
\label{fig:responsividade-wizard}
\legend{Fonte: capturas locais da aplicação, sem autenticação.}
\end{figure}
\FloatBarrier

A análise automatizada selecionou as regras A e AA disponíveis no axe para WCAG~2.0, 2.1 e 2.2. A Tabela~\ref{tab:verificacao-axe} resume as ocorrências reportadas por regra e cenário. Uma ocorrência corresponde a um elemento identificado pela ferramenta; as contagens de cenários distintos não devem ser somadas como problemas únicos. O relatório também contém itens inconclusivos, sobretudo de contraste, que requerem inspeção manual.

\begin{table}[htbp]
\centering
\caption{Ocorrências reportadas pelo axe no recorte do assistente.}
\label{tab:verificacao-axe}
\footnotesize
\renewcommand{\arraystretch}{1.2}
\begin{tabular}{|p{6.3cm}|c|c|c|c|}
\hline
\textbf{Regra} & \textbf{Claro} & \textbf{Claro} & \textbf{Escuro} & \textbf{Escuro} \\
 & 1440 & 390 & 1440 & 390 \\\hline
Botão sem nome acessível & 1 & 1 & 1 & 1 \\\hline
Contraste insuficiente & 14 & 14 & 5 & 5 \\\hline
Título de documento ausente & 1 & 1 & 1 & 1 \\\hline
Link sem nome acessível & 0 & 1 & 0 & 1 \\\hline
Elementos interativos aninhados & 6 & 6 & 6 & 6 \\\hline
Região rolável sem foco & 3 & 0 & 3 & 0 \\\hline
Alvo de interação pequeno & 1 & 1 & 1 & 1 \\\hline
\end{tabular}
\legend{Fonte: relatório automatizado da revisão de 12 de setembro de 2026.}
\end{table}

Esses achados impedem tratar a acessibilidade como requisito concluído. A aprovação da transição por teclado mostra apenas que os dois passos testados são acionáveis nesse contexto. Permanecem necessárias a correção e a inspeção dos problemas reportados, a verificação do foco visível e dos diálogos, a análise de outras páginas e testes com tecnologias assistivas, conforme a proposta da Seção~\ref{sec:protocolo-avaliacao}. O \textit{script} e o relatório estruturado da verificação acompanham este texto no repositório, nos diretórios \texttt{scripts} e \texttt{verificacao}.
